\documentclass[10pt]{article}

\usepackage[utf8]{inputenc}

\usepackage[T1]{fontenc}

\usepackage{amsmath}

\usepackage{amsfonts}

\usepackage{amssymb}

\usepackage[version=4]{mhchem}

\usepackage{stmaryrd}

\usepackage{CJKutf8}

\usepackage{mathrsfs}

\usepackage{graphicx}

\usepackage[export]{adjustbox}

\graphicspath{ {./images/} }



\begin{document}

рых случаях (напр., для полугрупп) понятие решеточной определяемости расширяют добавлением в заключение импликации условия «или $B$ антиизоморфна $A$ », так как антиизоморфные полугруппы также решеточно изоморфны. Классический пример решеточной определяемости доставляет первая основная теорема проективной геометрии (см. [1]), где в качестве $A$ рассматриваются векторные пространства над телами. Решеточно определяющимися являются также всякая абелева группа, содержащая два независимых элемента бесконечного порядка, всякая свободная группа (свободная полугруппа) и группа (полугруппа), нетривиально разложимая в свободное произведение, всякая нильпотентная группа без кручения, всякая коммутативная полугруппа с законом сокращения и без идемпотентов, всякая свободная полугруппа идемпотентов, свободная полурешетка более чем с двумя свободными образующими. При этом нередко оказывается, что каждое проектирование алгебры индуцируется ее изоморфизмом (или антиизоморфизмом). Класс однотипных алгебр \begin{CJK}{UTF8}{mj}气\end{CJK}̌ может содержать решеточно не определяющиеся алгебры, но обладать тем свойством, что из $A \in \mathscr{\mathcal { K }}$ и $\mathrm{Sub} B \cong \operatorname{Sub} A$ вытекает, что $B \in \mathscr{K}$; в этом случае говорят, что $\mathscr{K} \mathrm{p}$ еIII е т о ч н о о п р е де ля е т с я (или р е іI е то чн о $\quad$ а м к н у т); если при этом алгебры $B$ берутся только из класса $\mathcal{M}$, то к соответствующему термину добавляют «в классе $\mathscr{M}$ ». Среди решеточно замкнутых классов - класс всех разрепимых групп.



Многие ограничения, накладываемые на изучаемые алгебры, формулируются в терминах Р. п.; классический пример - условия минимальности и максимальности для подалгебр. Sub $A$ удовлетворяет условию максимальности тогда и только тогда, когда все подалгебры в $A$ конечно порожденные (см. также $\Gamma$ pyпna $c$ условием конечности, Полугруппа с условием конечности). В качестве других накладываемых на Р. п. ограничений рассматриваются такие теоретико-рететочные свойства как дистрибутивность, модулярность, различные виды полумодулярности, условие Жордана - Дедекинда, дополняемость, относительная дополняемость и т. Д. Напр., для группы $A$ Р. п. дистрибутивна тогда и только тогда, когда $A$ локально циклическая (т е о р е м а О р е); условия дистрибутивности изучены и в случае полугрупп, ассоциативных колец, модулей, алгебр Ли и др.



Наряду с изоморфизмами Р. п. рассматриваются д у ал и з м ы (т. е. антиизоморфизмы), гомоморфизмы. В случае, когда $A$ - топологическая алгебра, наиболее естественно сопоставлять ей решетку всех замкнутых подалгебр; соответствующая проблематика также активно разрабатывается.



Лum.: [1] Б ә р Р , Линейная алгебра и проективная геометрия, пер. с англ., М , 1955; [2] С у д з у к и М., Строение группы и строение структуры ее подгрупп, пер. с англ., М., 1960 [3] С к о р н я к о в Ј. А., Дедекиндовы структуры с дополнениями и регулярные кольца, М., $1961 ;[4]$ К о н П., Универсальниями и регулярные кольца, М., 1961; [4] К о н П., Универсаль-

ная алгебра, пер. с англ., М., 1968; [5] С а д о в с к и й Ј. Е., ная алгебра, пер. с англ., М., 1968; [5] С а д о в с к и й Ј. Е.,

"Успехи матем. наук», 1968, т. 23 , в. 3, с. $123-57 ;$ [6] А р ш ин о в М. Н., С а д о в с к и й Л. Е., там же, 1972, т. 27, в. 6, с. $139-80$; [7] II е в р и н Л. Н., “Сибирск. матем. ж.», 1962 , т. 3, № 3, с. 446-470; [8] Итоги науки Алгебра. 1964, М., 1966 , с. 237-74; [9] Итоги науки Алгебра Топология. Геометрия. 1968, М., 1970,с. 101-54; [10] Упорядоченные множества и решетки, в.' 3, Саратов, 1975, с. $50-74 . \quad$ J.Н.ІІеврин.



РЕШЕТКА С ДОПОЛНЕНИЯМИ - решетка $L$ с нулем 0 и единицей 1 , в к-рой для любого әлемента $a$ существует такой әлемент $b$ (наз. д о п о л н е н и е м э л е м е н т а $a$ ), что $a \vee b=1$ и $a \wedge b=0$. Произвольную решетку можно вложить в решетку, каждый элемент к-рой обладает единственным дополнением. Если для любых $a, b \in L$ интервал $[a, b]$ является Р. с д., то $L$ наз.



\includegraphics[max width=\textwidth, center]{2023_07_08_27d1522cfab57f574b03g-1}

н е н и я м и. Каждая модулярная Р. с д. является решеткой с относительными дополнениями. Решетка $L$ с нулем 0 называется: a) р е ші е т к о й с ч а с т и чны ми д о по лн е н и я м и, если каждый ее интервал вида [0,a], $a \in L$, является Р. с д.; б) р е ш е т к о й с о с л а бы м и до по лне ния м и, если для любых $a, b \in L$ существует такой элемент $c \in L$, что $a \wedge c=0$ и $b \wedge c \neq 0$; в) р е ше т к о й с по лу у о по лн ения м и, если для любого $a \in L, a \neq 1$, существует такой әлемент $b \in L, b \neq 0$, что $a \wedge b=0 ; r)$ р е ше т к о й с п с е в д о д о по лн е н и я м и, если для любого $a \in L$ существует такой элемент $a^{*}$, что $a \wedge x=0$ тогда и только тогда, когда $x \leqslant a^{*}$; д) р е ше т к о й с к в а з ид о п о л н е н и я м и, если для любого $x \in L$ существует такой элемент $y \in L, \quad y \neq x$, что $x \vee y$ является плотным элементом. Большую роль играют также репетки с ортодополнениями (см. Ортомодулярная решет$\kappa a)$. О связи между различными типами дополнений в решетках см. [4].



Лит.: [1] Б и р к г о ф Г., Теория структур, пер. с англ., М., 1952; [2] С к о р н я к о в Л. А., Элементы теории структур, нениями и регулярные кольца, М., 1961; [4] G r і 11 е t $\mathrm{P}$. А. V a r l e t J. C., "Bull. Soc. roy. Sci. Liège», 1967, t. 36, № 11-12, p. 628-42. же, что структурно упорядоченная группа.



РЕIIЕТЧАТОЕ РАСПРЕДЕЛЕНИЕ - дискретное вероятностное распределение, сосредоточенное на множестве точек вида $a+n h$, где $h>0, a$ - действительное число, $n=0, \pm 1, \pm 2, \ldots$. Число $h$ наз. ІІІ а г о м Р.p., и если ни при каких $a_{1}$ и $h_{1}>h$ распределение не сосредоточено на множестве вида $a_{1}+n h_{1}, n=0, \pm 1, \pm 2, \ldots$, то шаг $h$ наз. м а к с и м а л ь н ы м. Частным случаем Р. р. является арифметич. распределение.



Для того чтобы вероятностное распределение с характеристич. функцией $f(t)$ было р е іп е т ч а т ы м, необходимо и достаточно, чтобы существовало действйтельное число $t_{0} \neq 0$ такое, что $\left|f\left(t_{0}\right)\right|=1$; причем $h$ является максимальным шагом тогда и только тогда, когда $|f(t)|<1$ при $0<t<2 \pi / h$ и $|f(2 \pi / h)|=1$. Характеристич. функция Р. р. является периодич. функцией.



Формула обращения для Р. р. имеет вид



\[

p_{n}=\frac{h}{2 \pi} \int_{|t|<\pi / h} e^{-i t(a+n h)} f(t) d t,

\]



где $p_{n}$ - вероятность, к-рую Р. р. приписывают точке $a+n h, f(t)$ - соответствующая характеристич. функция. Справедливо также равенство



\[

\sum_{n=-\infty}^{\infty} p_{n=\frac{h}{2 \pi}}^{2} \int_{|t|<\pi / h}|f(t)|^{2} d t

\]



Свертка двух Р. p. с піагами $h_{1}$ и $h_{2}$ является P.p. тогда и только тогда, когда $h_{1} / h_{2}-$ рациональное число.



При исследовании предельного поведения сумм независимых случайных величин, имеющих Р. р., основной результат центральной предельной теоремы о сходимости к нормальному распределению существенно дополняется локальными теоремами для Р. р. Простейшим примером локальной теоремы для Р. р. является Лапласа теорема. Ее обобщением служит следующее утверждение: пусть $X_{1}, X_{2}, \ldots$ - последовательность независимых одинаково распределенных случайных величин c $\mathrm{E} X_{1}=m, \mathrm{D} X_{1}=\sigma^{2}$ и $S_{k}=X_{1}+\ldots+X_{k}$ и $X_{1}$ принимает значения вида $a+n h, h>0$; тогда если



\[

P_{k}(n)=\mathrm{P}\left\{S_{k}=k a+n h\right\},

\]



для выполнения при $k \rightarrow \infty$ предельного соотнопения $\left|\frac{\sigma \sqrt{k}}{h} P_{k}(n)-\frac{1}{\sqrt{2 \pi}} \exp \left\{-\frac{1}{2}\left(\frac{k a+n h-k m}{\sigma \sqrt{h}}\right)^{2}\right\}\right| \rightarrow 0$ равномерно относительно $n$, необходимо и достаточно, чтобы піаг $h$ был максимальным.



Лит.: [1] Г н е д е н к о Б. В., Курс теории вероятностей, 5 изд., М., 1969; [2] П е т р о в В. В., Суммы независимых случайных величин, М., 1972; [3] П р о х о р о в Ю. В., Р о з ан о в Ю. А., Теория вероятностей, 2 изд., М., 1973.





\end{document}